\section{Implementation}
\label{sec:impl}

The system is a small Python package (orchestrator, judge, tools,
cost, UI) built on LangGraph for the orchestrator state machine and
plain HTTP for the judge fan-out. All LLM traffic flows through
OpenRouter, which keeps the system model-agnostic and provides per-call
usage accounting that we prefer over static price tables.
\todo{Add codebase map figure or table.}

\begin{itemize}
  \item \textbf{Orchestrator runtime.} LangGraph drives the
        sequential tool-calling loop; state is the taxonomy plus a
        bounded scratchpad of past judge reports.
  \item \textbf{Judge fan-out.} Raw \texttt{requests.post} calls in a
        threadpool, sized to the cheap-model concurrency budget; no
        framework overhead per item.
  \item \textbf{Provider abstraction.} OpenRouter is the single
        backend, so swapping orchestrator or judge models is a CLI
        flag. Defaults: orchestrator
        \texttt{anthropic/claude-sonnet-4.6}, judge
        \texttt{meta-llama/llama-3.3-70b-instruct}.
  \item \textbf{Cost tracking.} Per-call cost is read from
        OpenRouter's native \texttt{usage.cost} field when present and
        falls back to a static price table otherwise; running totals
        are streamed to the UI.
  \item \textbf{Partial progress.} Taxonomy state is checkpointed to
        \texttt{taxonomy\_state.json} after every iteration and
        classifications are streamed line-by-line to
        \texttt{classifications.jsonl}, so the work done before a
        crash or kill is recoverable on disk; the UI surfaces the
        partial state for inspection. We do not currently support
        re-launching a killed run from the saved state, flagged in
        \S\ref{sec:limitations}.
  \item \textbf{Streamlit UI.} Live view of cost, taxonomy state,
        per-iteration revise ops, and per-item classifications, with
        configurable model dropdowns.
\end{itemize}
